\thispagestyle{empty}
\setstretch{1.15} % Zeilenspacing
\section*{Abstract}

\bigskip 
This thesis explores the use of artificial intelligence for anomaly detection in electric vehicle supply equipment (EVSE), with a particular focus on wallbox charging units in private environments. The objective is to reliably detect cyberattacks and abnormal behavior by analyzing hardware and kernel event data using machine learning techniques, specifically, Long Short-Term Memory (LSTM) networks.\\
After reviewing established approaches, both publicly available datasets (hereafter referred to as CICEVSE2024) and self-collected data are analyzed, with decision trees applied for feature selection before training the LSTM models.\\
The implemented LSTM-based detection system is evaluated on both datasets and achieves high detection accuracy, especially for flooding attacks. The results demonstrate that careful feature selection significantly reduces model complexity and resource consumption, enabling real-time deployment on embedded devices. Nonetheless, limitations remain, particularly regarding cross-platform transferability and the black-box nature of LSTM models. Future work will focus on broadening the coverage of attack scenarios and extending detection to distributed, multi-station setups. Another priority will be ensuring that the detection algorithm always operates within a reserved portion of system resources. This would guarantee that Prometheus queries and anomaly detection remain functional even under severe flooding attacks.\\
Overall, this research contributes to enhancing the security and robustness of monitoring solutions for critical electric vehicle charging infrastructure.